% Written by ZeelandCLD's build.py from its model: edit the model there, never this file.
% Columns: Combined loop, Loops, Mechanism
C1 Salt builds, and farms respond in steps & R1, B1, B2, B3 & \emph{Root zone salinity} damages the \emph{Soil structure}, so less rain soaks in to refill the lens: the \emph{Rainwater lens thickness in the low polders} falls over years and \emph{Root zone salinity} rises further (R1). As \emph{Harvest lost to salt} rises, and after the dry summers that bring the damage to notice, farmers respond in different ways, and at different speeds: \emph{Adoption of salt-tolerant varieties} of their current crops, with the same machines and storage (B1); \emph{Adoption of new crops}, new to the island's rotations, which need a market and their own machines (B2); and, over decades, less \emph{Arable land in the polders} where arable farming stops paying, which raises the \emph{Water board's plans to raise ditch levels}, since its levels follow the land's use (B3). The first two lower the \emph{Harvest lost to salt} and leave the salt in the soil; only the third lowers the salt. The order is a reading of the costs, since the first two wait alike for the damage to show, and of the decades B3 takes; and the links from \emph{Harvest lost to salt} to \emph{Adoption of salt-tolerant varieties} and from \emph{Agricultural land value} to \emph{Arable land in the polders} are hypotheses for the group model building session. \\
C2 Low ditch levels: workable fields now, more demand later & B4, R2 & A lower \emph{Ditch level (NAP)} gives a greater \emph{Polder drainage depth}, which raises the \emph{Spring workability of the fields} within the season, so \emph{Farmers' demand for deep polder drainage} eases (B4). The same lower level raises \emph{Saline seepage (kwel)}, which over years lowers the \emph{Rainwater lens thickness in the low polders}. \emph{Root zone salinity} and \emph{Harvest lost to salt} rise, \emph{Farm income} and \emph{Agricultural land value} fall. \emph{Financial leverage}, \emph{Banks' risk assessment of farm loans}, the \emph{Risk premium on farm loans} and \emph{Debt payments} rise, and thus the farm's \emph{Dependence on high-margin salt-sensitive crops}, which in spring needs deep drainage, grows: the demand returns higher (R2). The relief comes first and the extra demand years later, also known as the “Fixes that fail” archetype \parencite[p.~16]{kim-2000}. The link from \emph{Debt payments} to \emph{Dependence on high-margin salt-sensitive crops} is a hypothesis for the group model building session. \\
C3 Debt and capital together keep a farm in its crops & R2, R3 & \emph{Debt payments} raise the farm's \emph{Dependence on high-margin salt-sensitive crops} (R2, a hypothesis), and so does the \emph{Capital lost on switching to new crops} (R3, on sources). Each loop raises the variable that drives the other: more dependence brings more \emph{Spending on keeping the current crops}, so more \emph{Machines and storage for the current crops} and more capital a switch would write off (R3). This causes more \emph{Farmers' demand for deep polder drainage} (R2). Dependence on current crops holds back the adaptation of new crops. Salt, on its own, lowers that dependence over the years as the current crops yield worse harvests due to increasing root-zone salinity. \\
C4 The quick answer to a dry summer goes first & B5, B6, B7, B8 & Three loops are connected to \emph{Crop water shortfall}, which happen at different speeds: \emph{Groundwater extraction for irrigation} within the season (B5); \emph{Freshwater storage measures}, over the years the \emph{Rainwater lens thickness in the low polders} takes to grow (B6); and a \emph{Freshwater supply route to the island}, after the years it takes to decide, pay for and build a supply route, which the island does not have yet (B7; \cite[Regenwater section]{binnenlands-bestuur-2023-zoetwater}). Whatever eases the shortfall first lowers the pressure for the slower answers. Extraction is tied to what it draws on (B8): \emph{Pumping the water board allows} depends on the \emph{Freshwater lens thickness under creek ridges and dunes}, since the water board's rules let farms pump for irrigation without a permit, up to 60 m³ an hour, 3,000 m³ a quarter and 8,000 m³ a year and notified at least four weeks ahead, where its map shows a fresh groundwater body outside its vulnerable areas, one at least 15 m thick or reaching the geohydrological base \parencite[arts. 1.1 and 4.8 b and f]{scheldestromen-2026-waterschapsverordening}. \\
C5 Storing: more worthwhile, harder to pay for & R4, B6 & \emph{Freshwater storage measures} lower the \emph{Winter rain drained away}, so more rain refills the lens and, over years, the \emph{Rainwater lens thickness in the low polders} rises; \emph{Freshwater availability for crops} rises, the \emph{Crop water shortfall} falls and \emph{Crop yield} rises. There the two loops part. A larger \emph{Crop water shortfall} makes storing more worthwhile to a farmer (B6), and through a lower \emph{Crop yield} leaves less \emph{Farm income} and less \emph{Farmers' capacity to invest} (R4), so the same dry year pushes storage up and down. Farms with the means store and farms without may not; which pull is stronger differs per farm. The link from \emph{Farmers' capacity to invest} to \emph{Freshwater storage measures} is a hypothesis for the group model building session. \\
C6 Salt losses raise the wish to switch and cut the means & B2, R6, B9, B10 & A higher \emph{Harvest lost to salt} makes \emph{Adoption of new crops} more worthwhile, after a delay (B2). The same losses lower \emph{Farm income} and \emph{Agricultural land value}, so \emph{Financial leverage} and the \emph{Banks' risk assessment of farm loans} rise, \emph{Bank credit to farms} falls and the \emph{Farmers' capacity to invest} in the switch falls (R6); less credit keeps leverage down (B9). A field that moves saves harvest from salt (R6) and raises the \emph{Margin given up on high-margin crops} (B10). The level of salt in the soil of the farmland decides whether switching is economically beneficial or not. \\
C7 New crops stay stuck or take off & R7, B10 & With few hectares of \emph{Adoption of new crops} there is little harvest for a plant to process, no \emph{Local processing capacity for new crops}, and transport keeps the \emph{Net price for new crops after transport} low (R7). Every field that starts being used for new crops raises the \emph{Margin given up on high-margin crops}, which \emph{Farm income} and \emph{Bank credit to farms} must carry (B10): adoption stays low. With enough hectares a local plant becomes worthwhile, the net price rises and more farmers grow the crop. How many hectares that takes is unclear from the literature and thus not clear in the CLD. \\
C8 One ditch level between four goals & B4, B11, B12, B13 & Four balancing loops act on the \emph{Water board's plans to raise ditch levels}, each towards its own goal. Fields dry enough for heavy machinery in spring (\emph{Spring workability of the fields}, B4); the fresh water lenses the board wants to keep as a measure against \emph{Saline seepage (kwel)} (B11); the optimum Polder drainage depth the board's norms aim at for each use and soil type (B12); and a plan the objecting owners can accept (\emph{Objections to land-use plans}, B13). The goals cannot all be met at once, so the level settles where those pressures balance. How far seepage weighs against the other goals in the board's plans, cannot be determined from the literature found. \\
C9 Objections to one plan shape the next & B13, R8, R9 & \emph{Objections to land-use plans} delay, change or stop a plan (B13). Two things follow: The ditch levels keep up the \emph{Saline seepage (kwel)} and \emph{Root zone salinity} that lower \emph{Agricultural land value}, and a lower value may bring more objections (R8). Plans changed or dropped lower \emph{Farmers' trust in government}, which may bring more objections to the next plan (R9). Whether objections settle at a plan the owners accept or keep growing is a question for the group model building session: two of its links are hypotheses, from \emph{Agricultural land value} to \emph{Objections to land-use plans} and from \emph{Objections to land-use plans} to \emph{Farmers' trust in government}. \\
C10 Falling land value: levels held, or land leaving arable & R2, R8, B3 & The three loops share the chain from the \emph{Water board's plans to raise ditch levels} through \emph{Ditch level (NAP)}, \emph{Saline seepage (kwel)}, \emph{Rainwater lens thickness in the low polders}, \emph{Root zone salinity} and \emph{Harvest lost to salt} to \emph{Farm income} and \emph{Agricultural land value}, and part ways there. As land loses value, \emph{Financial leverage} and \emph{Debt payments} lead farmers to ask for deep drainage (R2), and owners may object harder to raises (R8, a hypothesis): both keep the ditch levels low and the salt coming. Where arable farming stops paying, \emph{Arable land in the polders} falls, the board has more reason to raise the levels, and the salt falls back (B3). The same fall in value sends the water board opposite signals; which is stronger decides whether the decline persists or stops. \\
